%% kapa: exam=1 round=20 year=2009 period=1 form=A subject=civil_law subject_ko=민법 seq=8 number=8
%% source: Q-Net 감정평가사 기출문제 (cst00309) artlSeq=5024178 file=extracted/1008/5024178/2009년 제20회 감정평가사 1차시험 기출문제/1교시.hwp
%% conversion: kordoc 4.24.51 -> tools/gen1.py (auto); review=proofread-sonnet; flags=none
\begin{question}{8}
甲이 도서관에서 자신의 노트북을 사용 중에 잠시 자리를 비운 사이에 乙이 그 노트북을 훔쳐서 이러한 사실을 모르는 丙에게 무상으로 빌려주었다. 다음 중 옳은 것은? (다툼이 있으면 판례에 의함)
\begin{choices}
  \item 甲은 점유침탈을 이유로 乙에게 노트북의 반환을 청구할 수 있으나 손해배상을 청구할 수는 없다.
  \item 丙은 포괄승계인에 해당하므로 甲은 丙에게 점유침탈을 이유로 노트북의 반환을 청구할 수 있다.
  \item 甲의 점유회수청구권은 침탈당한 날로부터 1년 내에 행사하여야 하는데, 이 기간은 출소(出訴)기간이다.
  \item 丙은 노트북에 대한 선의취득을 이유로 甲의 소유물반환청구를 거절할 수 있다.
  \item 甲은 乙에 대하여 소유권에 기하여 노트북의 반환을 청구할 수 있을 뿐 점유침탈을 이유로 노트북의 반환을 청구할 수 없다.
\end{choices}
\end{question}
\answer{3}
